% LaTeX-Präambel für TeX + Lilypond
% (LaTeX preamble for TeX + Lilypond - shared macros and packages of the whole book)
% @author: Hartmut Leister <hartmut.leister@gmail.com>

% ---------------------------------------------------------
% benötigte Variablen/Pakete - these have to be defined in the main .lytex file BEFORE
% including this preamble:
%
% \usepackage{ifthen} % wird für boolsche Flags benötigt
%
% % Titel zur Anzeige im Blattkopf
% \newcommand{\shortHead}{My title}
% % Datum zur Anzeige
% \newcommand{\theDate}{\today{}}
% % Author des Dokuments
% \newcommand{\theAuthor}{Hartmut Leister}
% % Logo-Dateien (nichts eintragen -> kein Logo)
% \newcommand{\logoSmallFileName}{logo.pdf}
% \newcommand{\logoLargeFileName}{logo.pdf}
% % Flag, ob Kopf- oder Fußzeilen angezeigt werden sollen (true, false)
% \newboolean{ShowHeader}
% \setboolean{ShowHeader}{false}
% ---------------------------------------------------------

\usepackage{ifxetex}
\ifxetex
    %xetex specific stuff
    \usepackage{xunicode,fontspec,xltxtra}
    \usepackage{relsize}
    \usepackage[ngerman]{babel}

    \setmainfont{GFS Didot}
    \setsansfont{Umpush}
\else
    % pdflatex specific configuration
    \usepackage[T1]{fontenc}
    \usepackage[utf8]{inputenc}
    \usepackage{mathptmx}%Times
    % Schriftartenfamile Linux Libertine + Biolinum
    \usepackage{libertine}
    \usepackage{german}
\fi

\usepackage{geometry}
% keep that in sync with scores/default.ly -> \paper
\geometry{  a5paper,
            left   = 1.0cm,
            right  = 1.0cm,
            top    = 2.0cm,
            bottom = 1.5cm,
            headheight = 25pt,
            heightrounded
            }
\pagestyle{empty}
\usepackage{graphicx}
\usepackage{wrapfig}

% symbols (e.g. \Cross)
\usepackage{marvosym}

% URLs and hyperrefs
\usepackage[breaklinks]{hyperref}
\usepackage{url}

\hypersetup{pdfborder=0 0 0}

\hypersetup{
%pdfpagemode = FullScreen,
pdfpagelayout = TwoPageRight,
baseurl = {},
pdftitle = {\shortHead},
pdfcreator = {TeX},
pdfproducer = {pdfTeX},
pdfauthor = {\theAuthor},
pdfsubject = {\shortHead},
pdfkeywords = {\theDate, \shortHead}
}

\usepackage{hyphenat}

\usepackage{mathtools}
% ---------------------------------------------------------
%  index generation
\usepackage{makeidx}
\makeindex
\renewcommand\indexname{Inhaltsverzeichnis}

% ---------------------------------------------------------
\usepackage{fancyhdr}
\ifShowHeader
    \pagestyle{fancy}
    \fancyhead{} % clear all header fields
    \fancyhead[LE]{\textsc{\shortHead}}
    % add small logo to header
    % height must be less than headheight above at geometry or we'll have mismatching heights
    \fancyhead[RO]{\includegraphics[height=20pt]{\logoSmallFileName}}
    \fancyfoot{} % clear all footer fields
    \fancyfoot[LE,RO]{\thepage}
    \fancyfoot[RE]{\theDate}
    \renewcommand{\headrulewidth}{0.4pt}
    \renewcommand{\footrulewidth}{0.4pt}
\fi

\setlength{\parindent}{0pt}
\setlength{\parskip}{8pt}

\usepackage{framed,color}
\definecolor{shadecolor}{rgb}{0.95,0.95,0.95}

% remove numberings from chapter, section, subsection
% (using just chapter* would not add them to the TOC
% (see https://tex.stackexchange.com/a/30202/19668)
\renewcommand{\thesection}{}
\renewcommand{\thechapter}{}

%\addsubsec
\makeatletter
\newcommand*\addsubsec{\secdef\@addsubsec\@saddsubsec}
\newcommand*{\@addsubsec}{}
\def\@addsubsec[#1]#2{\subsection*{#2}\addcontentsline{toc}{subsection}{#1}
  \if@twoside\ifx\@mkboth\markboth\markright{#1}\fi\fi
}
\newcommand*{\@saddsubsec}[1]{\subsection*{#1}\@mkboth{}{}}
\makeatother

% Aussehen von section, subsection
% bei Layout müssen drei Befehle: Schriftstärke, Serifen/keine, Größe
\makeatletter
\renewcommand\section{\@startsection
{section}{1}{0mm}%           % name, ebene, einzug
{1\baselineskip}%            % vor-abstand
{1\baselineskip}%            % nach-abstand
{\bfseries\sffamily\Large}%  % layout: fett, sans-serif, groß
}
\makeatother
\makeatletter
\renewcommand\subsection{\@startsection
{subsection}{2}{0mm}%        % name, ebene, einzug
{0\baselineskip}%            % vor-abstand
{0.1\baselineskip}%          % nach-abstand
{\bfseries\sffamily\large}%  % layout
}
\makeatother

% ---------------------------------------------------------
% song macros
\newcommand{\Lied}[1]
{
    % Layout für Lied (ohne Liednummer)
    \par
    \textbf{
        \large
        #1
    }
}
\newcommand{\LiedNummer}[2]{
    % Layout für Lied mit Angabe der Liednummer
    \Lied{#1}{\hfill (#2)}
}
